%%%%%%%%%%%%%%%%%%%%%%%%%%%%%%%%%%%%%%%%%%%%%%%%%%%%%%
% IDIOMA
%%%%%%%%%%%%%%%%%%%%%%%%%%%%%%%%%%%%%%%%%%%%%%%%%%%%%%

\usepackage[spanish]{babel}

%%%%%%%%%%%%%%%%%%%%%%%%%%%%%%%%%%%%%%%%%%%%%%%%%%%%%%
% GEOMETRÍA
%%%%%%%%%%%%%%%%%%%%%%%%%%%%%%%%%%%%%%%%%%%%%%%%%%%%%%

\usepackage[
top=0.65in,
bottom=0.65in,
left=0.65in,
right=0.65in
]{geometry}

%%%%%%%%%%%%%%%%%%%%%%%%%%%%%%%%%%%%%%%%%%%%%%%%%%%%%%
% CLASE BOOK
%%%%%%%%%%%%%%%%%%%%%%%%%%%%%%%%%%%%%%%%%%%%%%%%%%%%%%

\usepackage{fancyhdr}

\pagestyle{fancy}
\fancyhf{}
\fancyfoot[C]{\thepage}
\renewcommand{\headrulewidth}{0pt}

%%%%%%%%%%%%%%%%%%%%%%%%%%%%%%%%%%%%%%%%%%%%%%%%%%%%%%
% FUENTES
%%%%%%%%%%%%%%%%%%%%%%%%%%%%%%%%%%%%%%%%%%%%%%%%%%%%%%

\usepackage{fontspec}
\usepackage{microtype}

%%%%%%%%%%%%%%%%%%%%%%%%%%%%%%%%%%%%%%%%%%%%%%%%%%%%%%
% MATEMÁTICAS
%%%%%%%%%%%%%%%%%%%%%%%%%%%%%%%%%%%%%%%%%%%%%%%%%%%%%%

\usepackage{amsmath}
\usepackage{amssymb}
\usepackage{amsthm}
\usepackage{cancel}

%%%%%%%%%%%%%%%%%%%%%%%%%%%%%%%%%%%%%%%%%%%%%%%%%%%%%%
% GRÁFICOS Y TABLAS
%%%%%%%%%%%%%%%%%%%%%%%%%%%%%%%%%%%%%%%%%%%%%%%%%%%%%%

\usepackage{graphicx}
\usepackage{booktabs}
\usepackage{tabularx}
\usepackage{multirow}
\usepackage{multicol}
\usepackage{float}
\usepackage{pdflscape}
\usepackage{pgfplots}
\usepackage{subcaption}

\usetikzlibrary{arrows.meta}


\pgfplotsset{compat=1.18}
 
\definecolor{azul}{HTML}{1F5FAD}
\definecolor{naranja}{HTML}{D9730D}
\definecolor{verde}{HTML}{2E8B57}
\pgfplotsset{
  onda/.style={
    width=0.82\linewidth, height=4.2cm,
    axis lines=middle, axis line style={gray!70},
    xmin=-0.2, xmax=2*pi+0.35, ymin=-1.3, ymax=1.55,
    xtick={0,pi/4,pi/2,pi,3*pi/2,2*pi},
    xticklabels={$0$,$\frac{\pi}{4}$,$\frac{\pi}{2}$,$\pi$,$\frac{3\pi}{2}$,$2\pi$},
    ytick={-1,1},
    xlabel={$t$}, xlabel style={at={(ticklabel* cs:1)}, anchor=west},
    tick label style={font=\footnotesize},
    legend pos=outer north east,
    legend style={font=\footnotesize, draw=none, cells={anchor=west}},
    domain=0:2*pi, samples=200, clip=false,
  },
  base/.style={very thick, azul, no markers},
  imagen/.style={very thick, naranja, dashed, no markers},
  segunda/.style={very thick, verde, densely dotted, no markers},
  pico/.style={only marks, mark=*, mark size=1.8pt, draw=white, forget plot},
}
\tikzset{desplazamiento/.style={-latex, semithick, gray!80!black}}

%%%%%%%%%%%%%%%%%%%%%%%%%%%%%%%%%%%%%%%%%%%%%%%%%%%%%%
% COLORES
%%%%%%%%%%%%%%%%%%%%%%%%%%%%%%%%%%%%%%%%%%%%%%%%%%%%%%

\usepackage[dvipsnames,table]{xcolor}
\usepackage{colortbl}
%==================================================
% COLORES INSTITUCIONALES UNAL
%==================================================

\definecolor{verdeUN}{RGB}{148,180,59}      % Pantone 376C
\definecolor{rojoUN}{RGB}{166,28,49}        % Pantone 187C
\definecolor{rojoOscuroUN}{RGB}{118,35,47}  % Pantone 188C
\definecolor{verdeOscuroUN}{RGB}{70,107,63} % Pantone 7743C
\definecolor{grisOscuroUN}{RGB}{86,90,92}   % Pantone 425C
\definecolor{grisClaroUN}{RGB}{177,178,176} % Pantone 421C

%==================================================
% COLORES PARA CÓDIGO
%==================================================

\definecolor{bgcream}{HTML}{FFF8E7}
\definecolor{keywordcolor}{RGB}{70,107,63}
\definecolor{stringcolor}{RGB}{166,28,49}
\definecolor{commentcolor}{RGB}{86,90,92}

%%%%%%%%%%%%%%%%%%%%%%%%%%%%%%%%%%%%%%%%%%%%%%%%%%%%%%
% TIKZ
%%%%%%%%%%%%%%%%%%%%%%%%%%%%%%%%%%%%%%%%%%%%%%%%%%%%%%

\usepackage{tikz}

%%%%%%%%%%%%%%%%%%%%%%%%%%%%%%%%%%%%%%%%%%%%%%%%%%%%%%
% HIPERVÍNCULOS
%%%%%%%%%%%%%%%%%%%%%%%%%%%%%%%%%%%%%%%%%%%%%%%%%%%%%%

\usepackage{xurl}
\usepackage[hidelinks]{hyperref}

%%%%%%%%%%%%%%%%%%%%%%%%%%%%%%%%%%%%%%%%%%%%%%%%%%%%%%
% CAJAS
%%%%%%%%%%%%%%%%%%%%%%%%%%%%%%%%%%%%%%%%%%%%%%%%%%%%%%

\usepackage[most]{tcolorbox}

%%%%%%%%%%%%%%%%%%%%%%%%%%%%%%%%%%%%%%%%%%%%%%%%%%%%%%
% CAPTIONS Y LISTAS
%%%%%%%%%%%%%%%%%%%%%%%%%%%%%%%%%%%%%%%%%%%%%%%%%%%%%%

\usepackage{caption}
\usepackage{enumitem}

\setlength{\belowcaptionskip}{12pt}

%%%%%%%%%%%%%%%%%%%%%%%%%%%%%%%%%%%%%%%%%%%%%%%%%%%%%%
% NÚMEROS
%%%%%%%%%%%%%%%%%%%%%%%%%%%%%%%%%%%%%%%%%%%%%%%%%%%%%%

\usepackage{siunitx}

\sisetup{
table-number-alignment = center,
input-symbols = {+},
table-space-text-post = {+}
}

%%%%%%%%%%%%%%%%%%%%%%%%%%%%%%%%%%%%%%%%%%%%%%%%%%%%%%
% SECCIONES
%%%%%%%%%%%%%%%%%%%%%%%%%%%%%%%%%%%%%%%%%%%%%%%%%%%%%%

\usepackage{titlesec}

\titleformat{\section}
{\Large\bfseries}
{\thesection.}{0.5em}{}
[\vspace{0.3em}\color{rojoUN}\titlerule]


%%%%%%%%%%%%%%%%%%%%%%%%%%%%%%%%%%%%%%%%%%%%%%%%%%%%%%
% COLORBOX
%%%%%%%%%%%%%%%%%%%%%%%%%%%%%%%%%%%%%%%%%%%%%%%%%%%%%%

\newtcolorbox{cuadro_verde}[1]{
    enhanced,
    breakable,
    colback=white,
    colframe=verdeUN,
    colbacktitle=verdeUN,
    coltitle=white,
    fonttitle=\AncizarSansSemiBold,
    title=#1,
    boxrule=0.8pt,
    arc=2mm,
    left=3mm,
    right=3mm,
    top=2mm,
    bottom=2mm
}

\newtcolorbox{cuadro_verdeOscuro}[1]{
    enhanced,
    breakable,
    colback=white,
    colframe=verdeOscuroUN,
    colbacktitle=verdeOscuroUN,
    coltitle=white,
    fonttitle=\AncizarSansSemiBold,
    title=#1,
    boxrule=0.8pt,
    arc=2mm,
    left=3mm,
    right=3mm,
    top=2mm,
    bottom=2mm
}

\newtcolorbox{cuadro_rojo}[1]{
    enhanced,
    breakable,
    colback=white,
    colframe=rojoUN,
    colbacktitle=rojoUN,
    coltitle=white,
    fonttitle=\AncizarSansSemiBold,
    title=#1,
    boxrule=0.8pt,
    arc=2mm,
    left=3mm,
    right=3mm,
    top=2mm,
    bottom=2mm
}

\newtcolorbox{cuadro_rojo_Oscuro}[1]{
    enhanced,
    breakable,
    colback=white,
    colframe=rojoOscuroUN,
    colbacktitle=rojoOscuroUN,
    coltitle=white,
    fonttitle=\AncizarSansSemiBold,
    title=#1,
    boxrule=0.8pt,
    arc=2mm,
    left=3mm,
    right=3mm,
    top=2mm,
    bottom=2mm
}

\newtcolorbox{cuadro_grisClaro}[1]{
    enhanced,
    breakable,
    colback=white,
    colframe=grisClaroUN,
    colbacktitle=grisClaroUN,
    coltitle=white,
    fonttitle=\AncizarSansSemiBold,
    title=#1,
    boxrule=0.8pt,
    arc=2mm,
    left=3mm,
    right=3mm,
    top=2mm,
    bottom=2mm
}

\newtcolorbox{cuadro_grisOscuro}[1]{
    enhanced,
    breakable,
    colback=white,
    colframe=grisOscuroUN,
    colbacktitle=grisOscuroUN,
    coltitle=white,
    fonttitle=\AncizarSansSemiBold,
    title=#1,
    boxrule=0.8pt,
    arc=2mm,
    left=3mm,
    right=3mm,
    top=2mm,
    bottom=2mm
}


\newtcolorbox{nota}{
    enhanced,
    breakable,
    colback=white,
    colframe=verdeUN,
    borderline west={3pt}{0pt}{verdeUN},
    boxrule=0pt,
    arc=0mm,
    left=5mm,
    right=3mm,
    top=2mm,
    bottom=2mm
}


%%%%%%%%%%%%%%%%%%%%%%%%%%%%%%%%%%%%%%%%%%%%%%%%%%%%%%
% LISTINGS
%%%%%%%%%%%%%%%%%%%%%%%%%%%%%%%%%%%%%%%%%%%%%%%%%%%%%%

%%%%%%%%%%%%%%%%%%%%%%%%%%%%%%%%%%%%%%%%%%%%%%%%%%%%%%
% FUENTE JETBRAINS MONO
%%%%%%%%%%%%%%%%%%%%%%%%%%%%%%%%%%%%%%%%%%%%%%%%%%%%%%

%\setmonofont{Latin Modern Mono}
\setmonofont{IBM Plex Mono}
%\setmonofont{Source Code Pro}

%\setmonofont[
%Path=Config/Fuentes/JetBrains/,
%UprightFont=JetBrainsMono-Regular.ttf,
%BoldFont=JetBrainsMonoNL-Bold.ttf,
%ItalicFont=JetBrainsMono-Italic.ttf,
%BoldItalicFont=JetBrainsMonoNL-BoldItalic.ttf
%]{JetBrainsMono-Regular.ttf}


\usepackage{listings}

\lstdefinestyle{Rcode}{
language=R,
backgroundcolor=\color{bgcream},
keywordstyle=\color{keywordcolor}\bfseries,
commentstyle=\color{commentcolor}\itshape,
stringstyle=\color{stringcolor},
basicstyle=\fontsize{10}{14}\ttfamily\selectfont,
numbers=left,
numberstyle=\footnotesize\color{grisClaroUN},
numbersep=8pt,
frame=single,
rulecolor=\color{grisClaroUN},
breaklines=true,
showstringspaces=false,
morekeywords={library,ggplot,df2,zcounts},
xleftmargin=10pt,
framesep=6pt
}

\lstdefinestyle{Routput}{
basicstyle=\fontsize{10}{14}\ttfamily\selectfont,
backgroundcolor=\color{white},
frame=single,
rulecolor=\color{grisClaroUN},
breaklines=true,
showstringspaces=false,
numbers=none,
xleftmargin=10pt,
framesep=6pt
}

%%%%%%%%%%%%%%%%%%%%%%%%%%%%%%%%%%%%%%%%%%%%%%%%%%%%%%
% FUENTE PRINCIPAL DEL DOCUMENTO (SERIF)
%%%%%%%%%%%%%%%%%%%%%%%%%%%%%%%%%%%%%%%%%%%%%%%%%%%%%%

\setmainfont[
Path=Config/Fuentes/Serif/,
UprightFont=AncizarSerif-Regular.ttf,
ItalicFont=AncizarSerif-Italic.ttf,
BoldFont=AncizarSerif-Bold.ttf,
BoldItalicFont=AncizarSerif-BoldItalic.ttf
]
{AncizarSerif-Regular.ttf}

%%%%%%%%%%%%%%%%%%%%%%%%%%%%%%%%%%%%%%%%%%%%%%%%%%%%%%
% FUENTES SANS (PORTADA)
%%%%%%%%%%%%%%%%%%%%%%%%%%%%%%%%%%%%%%%%%%%%%%%%%%%%%%

\newfontfamily\AncizarSansMedium[
Path=Config/Fuentes/Sans/
]
{AncizarSans-Medium.ttf}

\newfontfamily\AncizarSansLight[
Path=Config/Fuentes/Sans/
]
{AncizarSans-Light.ttf}

\newfontfamily\AncizarSansLightItalic[
Path=Config/Fuentes/Sans/
]
{AncizarSans-LightItalic.ttf}

\newfontfamily\AncizarSansRegular[
Path=Config/Fuentes/Sans/
]
{AncizarSans-Regular.ttf}

\newfontfamily\AncizarSansSemiBold[
Path=Config/Fuentes/Sans/
]
{AncizarSans-SemiBold.ttf}

\newfontfamily\AncizarSansExtraBold[
Path=Config/Fuentes/Sans/
]
{AncizarSans-ExtraBold.ttf}

%%%%%%%%%%%%%%%%%%%%%%%%%%%%%%%%%%%%%%%%%%%%%%%%%%%%%%
% FUENTES SERIF AUXILIARES
%%%%%%%%%%%%%%%%%%%%%%%%%%%%%%%%%%%%%%%%%%%%%%%%%%%%%%

\newfontfamily\AncizarSerifMedium[
Path=Config/Fuentes/Serif/
]
{AncizarSerif-Medium.ttf}

\newfontfamily\AncizarSerifLight[
Path=Config/Fuentes/Serif/
]
{AncizarSerif-Light.ttf}

\newfontfamily\AncizarSerifLightItalic[
Path=Config/Fuentes/Serif/
]
{AncizarSerif-LightItalic.ttf}

\newfontfamily\AncizarSerifRegular[
Path=Config/Fuentes/Serif/
]
{AncizarSerif-Regular.ttf}

\newfontfamily\AncizarSerifSemiBold[
Path=Config/Fuentes/Serif/
]
{AncizarSerif-SemiBold.ttf}

\newfontfamily\AncizarSerifExtraBold[
Path=Config/Fuentes/Serif/
]
{AncizarSerif-ExtraBold.ttf}

%%%%%%%%%%%%%%%%%%%%%%%%%%%%%%%%%%%%%%%%%%%%%%%%%%%%%%
% TALLER 1: COLORES, NOTACIÓN, CAJAS Y GRÁFICOS
% Reemplaza el bloque que empieza en "% --- colores ---"
% (todo lo que se pegó del punto 1.1 al final del preámbulo).
% Va DESPUÉS de cargar xcolor, tcolorbox y pgfplots y de
% definir las fuentes Ancizar.
%%%%%%%%%%%%%%%%%%%%%%%%%%%%%%%%%%%%%%%%%%%%%%%%%%%%%%

% ------------------------------ colores -------------------------------
% Un color fijo por cada función; se repite en fórmulas, tabla y gráficos.
\definecolor{colCu}{HTML}{1F5FA8}   % cos t
\definecolor{colSu}{HTML}{D35400}   % sin t
\definecolor{colCd}{HTML}{1E8449}   % cos 2t
\definecolor{colSd}{HTML}{7D3C98}   % sin 2t
\definecolor{areaPos}{HTML}{AED6F1} % área por encima del eje
\definecolor{areaNeg}{HTML}{F5B7B1} % área por debajo del eje
\definecolor{ceroCol}{HTML}{B03A2E} % ceros de las funciones
\definecolor{resaltado}{HTML}{154360}

\newcommand{\Cu}{{\color{colCu}\cos t}}
\newcommand{\Su}{{\color{colSu}\sin t}}
\newcommand{\Cd}{{\color{colCd}\cos 2t}}
\newcommand{\Sd}{{\color{colSd}\sin 2t}}
\newcommand{\destacado}[1]{\boxed{\color{resaltado}#1}}

% ------------------------------ notación ------------------------------
% \ket lo define el paquete physics; esta línea solo actúa si no está cargado.
\providecommand{\ket}[1]{\lvert #1\rangle}
\newcommand{\bk}[2]{\langle #1\mid #2\rangle}   % producto hermítico <f|g>
\newcommand{\Lp}[2]{(#1,#2)_{L^2}}              % producto interno de L^2
\newcommand{\spanR}{\operatorname{span}_{\mathbb R}}
\newcommand{\spanC}{\operatorname{span}_{\mathbb C}}
\DeclareMathOperator{\Real}{Re}
\DeclareMathOperator{\Imag}{Im}

% ------------------------------- cajas --------------------------------
% Mismo aspecto que las cajas de los apuntes (cajaUN), pero con nombres y
% contador propios: así no chocan con "definicion" de los apuntes y la
% numeración del taller no depende de las secciones (Definición 1, 2, ...).
\tcbset{
  cajaTaller/.style={
    enhanced, breakable,
    colback=white, colframe=#1, colbacktitle=#1, coltitle=white,
    fonttitle=\AncizarSansSemiBold,
    boxrule=0.8pt, arc=2mm, left=3mm, right=3mm, top=2mm, bottom=2mm
  }
}
% Uso: \begin{definicionT}[Nombre opcional] ... \end{definicionT}
\newtcolorbox[auto counter]{definicionT}[1][]{
  cajaTaller=verdeOscuroUN,
  title={Definición~\thetcbcounter\ifx\relax#1\relax\else: #1\fi}
}
% Uso: \begin{lemaT}[Nombre opcional]{etiqueta} ... \end{lemaT}
% (la etiqueta permite citarlo con \ref{etiqueta})
\newtcolorbox[auto counter]{lemaT}[2][]{
  cajaTaller=rojoUN, label={#2},
  title={Lema~\thetcbcounter\ifx\relax#1\relax\else: #1\fi}
}
% Uso: \begin{observacionT} ... \end{observacionT}
\newtcolorbox{observacionT}{
  enhanced, breakable,
  colback=grisClaroUN!8, colframe=verdeOscuroUN,
  borderline west={3pt}{0pt}{verdeOscuroUN},
  boxrule=0pt, arc=0mm, left=5mm, right=3mm, top=2mm, bottom=2mm,
  before upper={{\AncizarSansSemiBold\color{verdeOscuroUN} Observación.}\ }
}

% ------------------------- estilo de gráficos -------------------------
% Se llama "ondaT" para no reemplazar el estilo "onda" de los apuntes.
\pgfplotsset{
  ondaT/.style={
    width=7.4cm, height=4.6cm,
    axis lines=middle,
    axis line style={-{Stealth[length=2mm]}, black!60},
    xmin=0, xmax=2*pi+0.6, ymin=-1.35, ymax=1.45,
    xtick={0,pi/2,pi,3*pi/2,2*pi},
    xticklabels={$0$,$\frac{\pi}{2}$,$\pi$,$\frac{3\pi}{2}$,$2\pi$},
    ytick={-1,1},
    tick label style={font=\footnotesize},
    xlabel={$t$},
    every axis x label/.style={at={(ticklabel* cs:1)},anchor=west},
    title style={font=\small, yshift=-1mm},
    samples=241, domain=0:2*pi,
    trig format plots=rad,
    clip=false,
  },
}
% Panel con áreas positivas (azul) y negativas (rojo) y ceros marcados.
% #1 expresión, #2 color de la curva, #3 título, #4 ceros, #5 extras
\newcommand{\panelarea}[5]{%
  \begin{tikzpicture}
  \begin{axis}[ondaT, title={#3}]
    \addplot[draw=none, fill=areaPos] {max(#1,0)} \closedcycle;
    \addplot[draw=none, fill=areaNeg] {min(#1,0)} \closedcycle;
    \addplot[very thick, #2] {#1};
    \addplot[only marks, mark=*, mark size=1.8pt, ceroCol] coordinates {#4};
    #5
  \end{axis}
  \end{tikzpicture}}